\documentclass[12pt]{article}

\usepackage[margin=1in]{geometry}
\usepackage{amsmath,amssymb,amsthm}
\usepackage{fancyhdr}
\usepackage{graphicx}
\usepackage{bm}
\usepackage{xcolor}
\usepackage{cancel}
\usepackage{hyperref}

%%%%%%%%%%%%%%%%%%%%%%
\pagestyle{fssancy}
\fancyhead[LO,L]{Yuval Kochavi}
\fancyhead[CO,C]{Radial Marshak BC at the Foam--Wall Interface}
\fancyhead[RO,R]{\today}
\fancyfoot[LO,L]{}
\fancyfoot[CO,C]{\thepage}
\fancyfoot[RO,R]{}
\renewcommand{\headrulewidth}{0.4pt}
\renewcommand{\footrulewidth}{0.4pt}
%%%%%%%%%%%%%%%%%%%%%%

\newcommand{\Ts}{T_s}
\newcommand{\Taxis}{T_{\mathrm{axis}}}
\newcommand{\Tint}{T_{\mathrm{int}}}
\newcommand{\Rcm}{R}
\newcommand{\xF}{x_F}
\newcommand{\sSB}{\sigma_{\mathrm{SB}}}
\newcommand{\lR}{\lambda_R}
\newcommand{\lext}{\ell_{\mathrm{ext}}}
\newcommand{\alphaW}{\alpha_w}
\newcommand{\Ew}{\dot{E}_w}
\newcommand{\pder}[2]{\frac{\partial #1}{\partial #2}}

\begin{document}

\section{Motivation}

In the current 2D cylindrical model, the radiation temperature field inside the
foam is decomposed into an axial Marshak wave solution and a radial eigenfunction.
The radial profile satisfies a diffusion eigenvalue problem with an effective
boundary condition at the foam--wall interface $r = \Rcm$.

The existing approach computes the wall albedo $\alphaW$ from the wall energy
loss parameterization and feeds it into the eigenvalue equation as a parameter.
However, $\alphaW$ itself depends on the interface temperature $\Tint$, which
in turn depends on the radial profile and hence on $\alphaW$---creating a
circular dependency.

Here we propose a \textbf{self-consistent} formulation: apply a proper Marshak
boundary condition at the radial interface $r = \Rcm$ using the known
temperature profile from $r = 0$ to $r = \Rcm$. This yields a Robin boundary
condition that simultaneously determines the radial eigenvalue $\kappa_0$,
the interface temperature $\Tint$, and the albedo $\alphaW$.


\section{Setup: Radial Temperature Profile}

\subsection{Radiation diffusion in cylindrical geometry}

Consider a cylindrical foam of radius $\Rcm$, with a radiation temperature
field $T(r, z, t)$. At a fixed axial position $z$ and time $t$, the radiation
energy density in the diffusion approximation satisfies:
\begin{equation}
    \frac{1}{r}\pder{}{r}\left(r\, D(r)\,\pder{U}{r}\right) = S(r)
    \label{eq:radial_diffusion}
\end{equation}
where $U = a T^4$ is the radiation energy density and $D = c/(3\kappa_R)$ is
the diffusion coefficient with $\kappa_R = 1/\lR$ the Rosseland mean opacity.

\subsection{Separation of variables}

We separate the radial dependence by writing:
\begin{equation}
    U(r, z, t) = U_0(z, t) \cdot \Phi(r)
    \label{eq:U_separation}
\end{equation}
where $U_0(z,t)$ is the on-axis ($r=0$) radiation energy density. The radial
eigenfunction $\Phi(r)$ satisfies:
\begin{equation}
    \frac{1}{r}\frac{d}{dr}\left(r\,\frac{d\Phi}{dr}\right) + \kappa_0^2\,\Phi = 0
    \label{eq:bessel_ode}
\end{equation}
with $\Phi(0) = 1$ (normalization) and $\Phi'(0) = 0$ (symmetry). The general
solution regular at $r = 0$ is:
\begin{equation}
    \boxed{\Phi(r) = J_0(\kappa_0\, r)}
    \label{eq:bessel_solution}
\end{equation}
where $J_0$ is the Bessel function of the first kind and $\kappa_0$ is the
radial eigenvalue, to be determined by the boundary condition at $r = \Rcm$.


\section{The Radial Marshak Boundary Condition}

\subsection{Half-range fluxes at the interface}

At the foam--wall interface $r = \Rcm$, we decompose the radiation field
into outgoing (foam $\to$ wall) and incoming (wall $\to$ foam) half-range
fluxes. In the $P_1$ (diffusion) approximation:
\begin{align}
    J^+(R) &= \frac{c}{4}\,U(R) + \frac{1}{2}\,F_{\mathrm{net}}(R)
    \qquad \text{(outgoing: foam $\to$ wall)}
    \label{eq:J_plus} \\[6pt]
    J^-(R) &= \frac{c}{4}\,U(R) - \frac{1}{2}\,F_{\mathrm{net}}(R)
    \qquad \text{(incoming: wall $\to$ foam)}
    \label{eq:J_minus}
\end{align}
where the net diffusive flux (positive outward) at the surface is:
\begin{equation}
    F_{\mathrm{net}}(R) = -D\,\pder{U}{r}\bigg|_{r=R}
    = -\frac{c}{3\kappa_R}\,\pder{U}{r}\bigg|_{r=R}
    \label{eq:F_net}
\end{equation}

Since $T$ decreases radially ($\partial U / \partial r < 0$), we have
$F_{\mathrm{net}}(R) > 0$: net energy flows outward into the wall.

\subsection{Wall reflection condition}

The wall absorbs a fraction $(1 - \alphaW)$ of the incident flux and reflects
(re-emits) a fraction $\alphaW$:
\begin{equation}
    J^-(R) = \alphaW\,J^+(R)
    \label{eq:wall_reflection}
\end{equation}

\subsection{Combining the conditions}

Substituting Eqs.~\eqref{eq:J_plus}--\eqref{eq:J_minus} into
\eqref{eq:wall_reflection}:
\begin{equation}
    \frac{c}{4}\,U(R) - \frac{1}{2}\,F_{\mathrm{net}}(R)
    \;=\; \alphaW\left[\frac{c}{4}\,U(R) + \frac{1}{2}\,F_{\mathrm{net}}(R)\right]
    \label{eq:combined_BC}
\end{equation}

Rearranging:
\begin{equation}
    \frac{c}{4}\,U(R)\,(1 - \alphaW)
    \;=\; \frac{1}{2}\,F_{\mathrm{net}}(R)\,(1 + \alphaW)
\end{equation}

\begin{equation}
    \boxed{U(R) = \frac{2(1 + \alphaW)}{c\,(1 - \alphaW)}\;F_{\mathrm{net}}(R)}
    \label{eq:robin_energy}
\end{equation}

Substituting $F_{\mathrm{net}} = -\frac{c}{3\kappa_R}\frac{\partial U}{\partial r}\big|_R$:
\begin{equation}
    U(R) = -\frac{2(1 + \alphaW)}{3\kappa_R\,(1 - \alphaW)}\;\pder{U}{r}\bigg|_{r=R}
    \label{eq:robin_U}
\end{equation}

This is a \textbf{Robin (mixed) boundary condition}. Defining the
\emph{extrapolation length}:
\begin{equation}
    \boxed{\lext = \frac{2(1 + \alphaW)}{3\kappa_R\,(1 - \alphaW)}}
    \label{eq:ell_ext}
\end{equation}
the boundary condition becomes:
\begin{equation}
    \Phi(R) = -\lext\;\Phi'(R)
    \label{eq:robin_phi}
\end{equation}


\section{Modified Eigenvalue Equation}

\subsection{Substituting the Bessel solution}

With $\Phi(r) = J_0(\kappa_0 r)$ and $\Phi'(r) = -\kappa_0 J_1(\kappa_0 r)$,
the Robin BC \eqref{eq:robin_phi} becomes:
\begin{equation}
    J_0(\kappa_0 R) = \lext\,\kappa_0\,J_1(\kappa_0 R)
\end{equation}

Substituting $\lext$ from \eqref{eq:ell_ext}:
\begin{equation}
    \boxed{J_0(\kappa_0 R) = \frac{2(1 + \alphaW)}{3\kappa_R\,(1 - \alphaW)}
    \;\kappa_0\;J_1(\kappa_0 R)}
    \label{eq:eigenvalue_full}
\end{equation}

Defining the dimensionless variable $x = \kappa_0 R$:
\begin{equation}
    J_0(x) = \frac{2(1 + \alphaW)}{3\kappa_R R\,(1 - \alphaW)}\;x\;J_1(x)
    \label{eq:eigenvalue_x}
\end{equation}

\subsection{Connection to the existing eigenvalue equation}

The existing code in \texttt{eigen\_bessel\_solver.py} solves:
\begin{equation}
    x\,J_1(x) = \varepsilon\,J_0(x)
    \label{eq:existing_eigen}
\end{equation}
with:
\begin{equation}
    \varepsilon = \frac{3}{4}\,(1 - \alphaW)\,\frac{\Rcm}{\lR}
    = \frac{3}{4}\,(1 - \alphaW)\,\kappa_R\,\Rcm
    \label{eq:existing_epsilon}
\end{equation}

Rewriting \eqref{eq:eigenvalue_x}:
\begin{equation}
    x\,J_1(x) = \frac{3\kappa_R R\,(1 - \alphaW)}{2(1 + \alphaW)}\;J_0(x)
\end{equation}

So the eigenvalue equation is equivalent to \eqref{eq:existing_eigen} with
$\varepsilon$ \emph{redefined} as:
\begin{equation}
    \boxed{\varepsilon_{\mathrm{new}} = \frac{3\kappa_R R\,(1 - \alphaW)}{2(1 + \alphaW)}
    = \frac{3}{2}\,\frac{1 - \alphaW}{1 + \alphaW}\,\frac{R}{\lR}}
    \label{eq:epsilon_new}
\end{equation}

\textbf{Comparison with the current formula:}
\begin{align}
    \varepsilon_{\mathrm{current}} &= \frac{3}{4}\,(1 - \alphaW)\,\frac{R}{\lR}
    \label{eq:eps_current} \\[6pt]
    \varepsilon_{\mathrm{new}} &= \frac{3}{2}\,\frac{1 - \alphaW}{1 + \alphaW}\,\frac{R}{\lR}
    \label{eq:eps_new}
\end{align}

For high albedo ($\alphaW \to 1$), both $\varepsilon \to 0$, giving $\kappa_0 \to 0$
and a flat radial profile (no lateral losses). For $\alphaW = 0$ (perfect absorber):
\begin{equation}
    \varepsilon_{\mathrm{current}}\big|_{\alphaW=0} = \frac{3}{4}\,\frac{R}{\lR}
    \qquad\text{vs.}\qquad
    \varepsilon_{\mathrm{new}}\big|_{\alphaW=0} = \frac{3}{2}\,\frac{R}{\lR}
\end{equation}
The new formula gives $2\times$ larger $\varepsilon$ for $\alphaW = 0$, which is the
correct Marshak extrapolation length $\lext = 2/(3\kappa_R)$.


\section{Limiting Behaviour}

\begin{itemize}
    \item \textbf{Perfect reflector ($\alphaW = 1$):} $\lext \to \infty$, so
    $\varepsilon \to 0$. The eigenvalue $\kappa_0 \to 0$, giving $\Phi(r) = J_0(0) = 1$
    everywhere---a flat radial profile. The wall reflects all incident radiation,
    so there is no radial loss. \checkmark

    \item \textbf{Perfect absorber ($\alphaW = 0$):} $\lext = 2/(3\kappa_R)$,
    recovering the standard Marshak extrapolation length for a vacuum boundary.

    \item \textbf{Optically thick foam ($\kappa_R R \gg 1$):} $\varepsilon \gg 1$
    and $\kappa_0 R \to j_{0,1} \approx 2.4048$ (the first zero of $J_0$), recovering
    the Dirichlet condition $\Phi(R) = 0$.

    \item \textbf{Optically thin foam ($\kappa_R R \ll 1$):} $\varepsilon \ll 1$
    and $\kappa_0 R \ll 1$, giving an approximately flat profile with a small
    correction $\Phi(r) \approx 1 - \kappa_0^2 r^2/4$.
\end{itemize}


\section{Self-Consistent Interface Temperature}

\subsection{The key idea}

Once the eigenvalue $\kappa_0$ is determined from \eqref{eq:eigenvalue_full},
the interface temperature follows directly from the radial profile:
\begin{equation}
    \boxed{T_{\mathrm{int}}^4 = T_{\mathrm{axis}}^4 \cdot J_0(\kappa_0 R)}
    \label{eq:T_interface}
\end{equation}
where $T_{\mathrm{axis}} = T_{\mathrm{axis}}(z,t)$ is the on-axis temperature from
the 1D Marshak march.

\subsection{Self-consistent albedo from the radial flux}

The net outward radiation flux at the interface is:
\begin{align}
    F_{\mathrm{net}}(R) &= -\frac{c}{3\kappa_R}\;\pder{(aT^4)}{r}\bigg|_{r=R}
    \nonumber \\
    &= \frac{c\,a}{3\kappa_R}\;\kappa_0\;T_{\mathrm{axis}}^4\;J_1(\kappa_0 R)
    \label{eq:F_at_R}
\end{align}

The power absorbed by the wall per unit axial length at position $z$ is:
\begin{equation}
    \dot{E}_w(z) = 2\pi R\,(1 - \alphaW)\;J^+(R)
    = 2\pi R\,\left[\frac{c}{4}\,a\,T_{\mathrm{int}}^4
    + \frac{1}{2}\,F_{\mathrm{net}}(R)\right](1 - \alphaW)
    \label{eq:Ew_from_flux}
\end{equation}

\textbf{Alternatively}, one can compute the absorbed flux purely from the
foam-side quantities, without needing the wall loss parameterization:
\begin{equation}
    \Ew = (1 - \alphaW)\;J^+(R) = J^+(R) - J^-(R) = F_{\mathrm{net}}(R)
    \label{eq:Ew_equals_Fnet}
\end{equation}
\textbf{This is a key result}: the net radial flux at the interface equals the
wall energy absorption rate per unit area. Since we can compute $F_{\mathrm{net}}(R)$
from the known Bessel profile, we can bypass the wall loss parameterization
$E_w(T_0, t)$ entirely for the albedo calculation.

\subsection{Self-consistency loop}

The procedure is:
\begin{enumerate}
    \item \textbf{Start} with an initial guess for $\alphaW$ (e.g., from the
    previous timestep, or from the wall loss parameterization).

    \item \textbf{Compute} $\varepsilon_{\mathrm{new}}$ from
    Eq.~\eqref{eq:epsilon_new} using the current $\alphaW$, $\kappa_R$, and $R$.

    \item \textbf{Solve} the eigenvalue equation \eqref{eq:existing_eigen}
    for $\kappa_0$ (already implemented in \texttt{eigen\_bessel\_solver.py}).

    \item \textbf{Extract} the interface temperature from
    Eq.~\eqref{eq:T_interface}: $T_{\mathrm{int}} = T_{\mathrm{axis}}\,[J_0(\kappa_0 R)]^{1/4}$.

    \item \textbf{Compute} the wall energy loss $\Ew$ using the wall loss
    parameterization evaluated at $T_{\mathrm{int}}$:
    \begin{equation}
        \Ew(z,t) = \frac{dE_w}{dt}\bigg|_{T_0 = T_{\mathrm{int}},\; t_{\mathrm{exp}}}
    \end{equation}
    (e.g., $E_w^{\mathrm{Au}} = 0.59\,T_0^{3.35}\,t_{\mathrm{ns}}^{0.59}$
    differentiated with respect to $t$).

    \item \textbf{Update} the albedo using the Rosen formula:
    \begin{equation}
        \alphaW^{\mathrm{new}} = 1 - \frac{\Ew}{\sSB\,T_{\mathrm{int}}^4}
        \label{eq:albedo_update}
    \end{equation}

    \item \textbf{Iterate} steps 2--6 until $|\alphaW^{\mathrm{new}} - \alphaW^{\mathrm{old}}| < \delta$
    (typically 2--3 iterations suffice).
\end{enumerate}


\section{Physical Interpretation}

\subsection{Why this is better}

The current approach has a conceptual issue: the albedo $\alphaW$ is computed
from the 1D surface temperature $\Ts$ (the Marshak march solution), but this
temperature corresponds to the \emph{on-axis} value at $r = 0$. The actual
temperature at the foam--wall interface is lower due to the radial profile:
\begin{equation}
    T_{\mathrm{int}} = T_{\mathrm{axis}}\cdot[J_0(\kappa_0 R)]^{1/4} < T_{\mathrm{axis}}
\end{equation}

Since the wall loss scales as a high power of temperature
($\Ew \propto T^{3.35}$ for gold), using $T_{\mathrm{axis}}$ instead of
$T_{\mathrm{int}}$ \textbf{overestimates} the wall energy loss, which in turn
\textbf{underestimates} the albedo.

The self-consistent approach corrects this by:
\begin{enumerate}
    \item Using the physically correct interface temperature.
    \item Coupling the albedo and radial profile consistently through the
    Robin BC.
    \item Naturally recovering the correct limiting behaviours (perfect
    reflector/absorber, optically thick/thin).
\end{enumerate}

\subsection{Geometric picture: extrapolated endpoint}

The Robin BC \eqref{eq:robin_phi} is equivalent to saying that the
eigenfunction $\Phi(r) = J_0(\kappa_0 r)$ extrapolates to zero not at $r = R$
but at an \emph{effective radius}:
\begin{equation}
    R_{\mathrm{eff}} = R + \lext
    = R + \frac{2(1 + \alphaW)}{3\kappa_R(1 - \alphaW)}
\end{equation}

\begin{itemize}
    \item For $\alphaW = 0$: $R_{\mathrm{eff}} = R + \frac{2}{3\kappa_R}$
    (standard diffusion extrapolation length).
    \item For $\alphaW \to 1$: $R_{\mathrm{eff}} \to \infty$
    (flat profile, no effective boundary).
\end{itemize}

This is the cylindrical analogue of the classical ``extrapolated endpoint''
in planar slab diffusion theory.


\section{Summary of Key Equations}

\begin{enumerate}
    \item \textbf{Radial profile:} $T^4(r,z,t) = T_{\mathrm{axis}}^4(z,t)\cdot J_0(\kappa_0 r)$

    \item \textbf{Eigenvalue equation (unchanged form):}
    \begin{equation*}
        x\,J_1(x) = \varepsilon\,J_0(x), \qquad x = \kappa_0 R
    \end{equation*}

    \item \textbf{Modified $\varepsilon$ from the radial Marshak BC:}
    \begin{equation*}
        \varepsilon = \frac{3}{2}\,\frac{1 - \alphaW}{1 + \alphaW}\,\frac{R}{\lR}
    \end{equation*}

    \item \textbf{Interface temperature:}
    \begin{equation*}
        T_{\mathrm{int}} = T_{\mathrm{axis}}\cdot\left[J_0(\kappa_0 R)\right]^{1/4}
    \end{equation*}

    \item \textbf{Self-consistent albedo (Rosen formula):}
    \begin{equation*}
        \alphaW = 1 - \frac{\Ew(T_{\mathrm{int}},\,t_{\mathrm{exp}})}
        {\sSB\,T_{\mathrm{int}}^4}
    \end{equation*}
\end{enumerate}

\end{document}
